\chapter{As long as you remember her, you are not alone (The Different Story)}
Symmetrically enough, Puella Magi Madoka Magica: The Different Story (by Hanokage based on a story by Magica Quartet) focuses on the last remaining of the three key tenets that spinoffs have to deal with, unfairness. Introducing no new characters, it instead takes a fairly narrow "what if" approach to the series.

The first volume of the series is largely a character study of Mami and Kyoko, serving as a prequel to both the anime and the rest of the manga in which Mami takes Kyoko, a newly minted magical girl, on as her apprentice. After the death of Kyoko's family, she and Mami have a falling out due to Kyoko rejecting Mami's philosophy of hunting both witches and familiars to protect others. The second and third volumes skip forward to a timeline, not shown in the series, in which Mami has taken on Sayaka as a new apprentice magical girl. In this timeline, Sayaka contracted with Kyubey slightly earlier and was thus able to rescue Mami from being killed by Charlotte. Kyoko returns, and she, Mami, and Sayaka clash over the proper way to live as a magical girl, while Sayaka descends into depression.

The comic focuses repeatedly on the essential unfairness of the magical girl system, but in a different way than Madoka Magica. The TV series primarily focuses on the unfair way in which Kyubey withholds key information from the contracting girls, such as the transformation of their bodies or the true nature of witches, but the former never comes up in The Different Story, while the latter is present but only briefly remarked upon. Instead, the series primarily focuses on the isolation imposed on the girls, and the effect this has on their psyches. Of course this isolation and the resulting depression was a major focus of the Sayaka arc of Madoka Magica, but in the TV series the isolation was at least partially self-imposed, as Sayaka deliberately pushed others away after learning how her body had been transformed. In The Different Story, isolation is instead inevitable, a built-in feature of the magical girl system.

Mami fights throughout the manga to try to find a way out of her isolation. She is desperate for friendship, but the endless work of being both a student and a magical girl means she has no time to spend with classmates. Her efforts to reach out to other magical girls are rebuffed because she "foolishly" fights familiars instead of focusing her efforts solely on witches and the resulting Grief Seeds. She does manage to befriend Kyoko, but loses her to the same philosophy espoused by the other magical girls. Later she manages to bond with Madoka and Sayaka, but then Sayaka becomes a witch in part because she pursued Mami's policy of hunting familiars. In the end, even after Madoka gives her a way out—fighting as a magical girl in her stead, allowing Mami to rest—Mami still feels alone and unworthy, to the point of suicide.

Mami's example is the most prominent, but the same holds for the other two girls who descend into despair over the course of The Different Story. Kyoko isolates herself almost completely after the death of her family, rebuffing overtures of friendship from everyone except Madoka. Sayaka, meanwhile, deliberately pushes away Mami, Kyoko, and Kyousuke in her shame at having abandoned Hitomi to a witch.

Arguably, all of this isolation is self-inflicted, but at the same time it is depicted as a consequence of the unfairness of the world in general and Kyubey's magical girl system in particular. The central argument between Mami and other magical girls throughout the manga is over whether the quest for Grief Seeds has to take precedence over other concerns, such as protecting others; further, Kyubey outright states that Sayaka becoming a witch is a direct consequence of the over-concentration of magical girls in Mitakihara, such that it becomes inevitable the weakest and least experienced will not be able to gain enough Grief Seeds. Mami's inability to hold together even a two-person team is no accident; it is a feature of a system designed to maximize the despair of magical girls in order to accelerate their transformation into witches. Two magical girls will inevitably compete over Grief Seeds, if something else doesn't tear them apart first.

Kyoko's story projects that unfairness into a wider world. Her father, like Kyubey, equates magical girls and witches, but he takes it further: he sees no difference between the two. He rejects Kyoko utterly as soon as he learns what she has become, claiming that she has deceived him and corrupted his teachings. He cannot conceive of magic or feminine power as anything other than witchcraft, and refuses to listen to Kyoko's explanations, leading him to turn to drink, then abuse, and finally murder-suicide.

That the magical girls are not to blame for their isolation is emphasized by the common thread of erroneous guilt through all three. Mami blames herself for her family's death because she didn't wish to save them, even though it's hardly unreasonable for a dying person to be focused on the wish to live to the exclusion of all else. In contrast, Kyoko blames herself for her own family's death because she wished to save them, even though she was helping them, and it was only her father's choice to refuse to listen or try to understand that killed them. This kind of guilt is relatively common in people who experience grief, loss, and trauma, and is a major symptom of post-traumatic stress disorder \cite{ref106}. Sayaka as well blames herself for not saving Hitomi from a witch, even though, as Madoka points out, she wasn't certain it was Hitomi at the time and had no way of knowing there was a witch involved.

All three girls isolate themselves out of guilt. Sayaka and Kyoko do so more obviously, by verbally (and sometimes physically) attacking those who try to befriend them. Mami also isolates herself, however, as she reveals in the final issue: she creates a false, "perfect"-seeming self in order to keep others from seeing her needs. This, as the reader knows from having seen Kyoko's perspective, only serves to heighten those others' feelings of inadequacy and isolation, perpetuating the cycle of loneliness and despair. Such self-isolation is another symptom of post-traumatic stress disorder \cite{ref107}, and the implication is thus that Kyubey's system is inherently traumatizing, subjecting every girl who passes through it to combat stress coupled with painful isolation.

Ultimately, Madoka finds a solution, wishing Sayaka back to life and taking Mami's place as a magical girl, but it is not enough. As the final sequence, which appears to be Mami's dying dream after her suicide, shows, only a world where all five magical girls—Mami, Kyoko, Sayaka, Madoka, and Homura—are able to work together and be friends with one another, and where Mami has a family as well, is enough to satisfy Mami's loneliness.

Unfortunately, such a world is not without issues of its own, as we will see in discussing Rebellion.
